\section{Zagrożenia bezpieczeństwa}
Wraz z rozwojem dużych modeli językowych pojawiają się ich nowe, groźne podatności. Udoskonalaniu mechanizmów bezpieczeństwa towarzyszy również rozwój coraz bardziej zaawansowanych metod ataku.
Stanowi to istotne wyzwanie dla twórców systemów AI, którzy dążą do rozbudowy funkcjonalności modeli jednocześnie minimalizując ryzyko ich nadużycia.
Potężne możliwości generowania tekstu przez duże modele językowe niosą ze sobą również zagrożenie takie jak tworzenie szkodliwych treści, zarówno w sposób nieintencjonalny jak i w wyniku przeprowadzonego ataku~\cite{zhang2025guardiansoffenderssurveyharmful}.
Generowane w ten sposób niedozwolone odpowiedzi można uporządkować w cztery kategorie~\cite{zhang2025guardiansoffenderssurveyharmful}.
Jedną z nich może być toksyczność obejmująca wulgaryzmy, wyzwiska czy uprzedzenia kulturowe. Drugą kategorią jest nękanie, którego celem jest zastraszanie czy kierowanie gróźb. Kolejną kategorią jest mowa nienawiści. Ostatnią kategorię tworzą treści obarczone uprzedzeniami, które skierowane są do osób o określonej rasie, płci, religii czy orientacji.
Rozwój bezpieczeństwa dużych modeli językowych jest kluczowy, by niwelować możliwość generowania takich treści i dostarczania ich użytkownikom.

Według organizacji OWASP najistotniejszym zagrożeniem dla dużych modeli językowych są ataki typu wstrzykiwanie poleceń (ang. prompt injection)~\cite{owasp_llm_top10}.
Podatność ta występuje, gdy dane wejściowe przekazane przez użytkownika modyfikują zachowanie lub odpowiedź modelu w sposób nieprzewidziany przez twórców.
Skutki udanego ataku mogą wykraczać poza samo generowanie szkodliwych treści
Mogą prowadzić do zakłóceń prawidłowego działania modelu, blokowania funkcjonalności, skłonienia systemu do odbiegania od pożądanego zachowania oraz wycieku prywatnych danych przetwarzanych przez LLM~\cite{nist_taxonomy_2025}.
Zagrożeniem wynikającym z tego typu ataków może być również obejście mechanizmów obronnych i zasad etycznych wbudowanych przez twórców modelu, przejęcie kontroli nad wywoływaniem zewnętrznych narzędzi, czy również wydobycie ukrytych informacji z kontekstu rozmowy, takich jak przykładowo instrukcja systemowa (ang. system prompt)~\cite{nist_taxonomy_2025}.